% --- SEZIONE 2: MODELLO DI MEMORIA E STACK ---
\section{Ottimizzazione e Strutture Dati}

\subsubsection{Switch Case}
\begin{itemize}
    \item \textbf{Pro:} Numero di cicli di clock definito e prevedibile.
    \item \textbf{Contro:} La condizione deve essere necessariamente una costante di tipo \texttt{int}.
    \item \textbf{Nota:} L'ordine dei \texttt{case} è rilevante per l'efficienza della ricerca.
\end{itemize}

\subsubsection{Efficienza delle Strutture Dati}
\begin{itemize}
    \item \textbf{Array Locali:} Molto \textbf{inefficienti}. 
    alto numero di cicli solo per l'inizializzazione nello stack frame 
    \item \textbf{Strutture:} È consigliabile \textbf{definire le strutture globalmente}. L'inizializzazione locale su stack comporta un overhead elevato di cicli di clock.
    \item \textbf{Bitfields:} Estremamente \textbf{inefficienti}. 
    Sebbene risparmino memoria, richiedono complesse operazioni di
     \textit{Read-Modify-Write} (LDR, Shift, Mask, OR, STR)
\end{itemize}